\originalchapter{相对论运动学与相空间}\label{ch:kinematics}
\originalsection{自然单位与不变量}
粒子物理分编取 $\hbar=c=1$, 四矢量 $p^\mu=(E,\mathbf p)$, 度规 $g_{\mu\nu}=\operatorname{diag}(1,-1,-1,-1)$. 壳上关系为 $p^2=m^2$. 恢复单位时质量写 MeV/$c^2$, 动量写 MeV/$c$, 截面乘 $(\hbar c)^2$. 若计算给 $1\,\mathrm{GeV}^{-2}$, 对应约 $0.3894\,\mathrm{mb}$.

沿 $z$ 的主动 boost 为 $E=\gamma(E^*+\beta p_z^*)$, $p_z=\gamma(p_z^*+\beta E^*)$, 横动量不变. 单粒子能量和动量随参考系变化, 四动量总和的平方则不变. 对 $a+b\to c+d$, 定义 $s=(p_a+p_b)^2$, $t=(p_a-p_c)^2$, $u=(p_a-p_d)^2$; 有 $s+t+u=\sum_i m_i^2$. 正的 $s$ 是质心能量平方, 不能写成四倍束流能量而漏掉平方.

对静止靶 $s=m_a^2+m_b^2+2m_bE_a$, 对等能反向相同质量束流 $s=4E^2$. 高能固定靶的 $\sqrt s$ 只按 $\sqrt{E_a}$ 增长, 对撞束流则按 $E$ 增长. 这解释为什么质心能量与实验室能量不是同一指标.

\originalsection{两体衰变}
$A\to B+C$ 在 $A$ 静止系满足 $\mathbf p_B=-\mathbf p_C$. 由 $p_C=p_A-p_B$, $m_C^2=m_A^2+m_B^2-2m_AE_B$, 得
\begin{equation}\label{eq:two-body}
 E_B=\frac{m_A^2+m_B^2-m_C^2}{2m_A},\quad
 p=\frac{\sqrt{[m_A^2-(m_B+m_C)^2][m_A^2-(m_B-m_C)^2]}}{2m_A}.
\end{equation}
第一括号给阈值条件 $m_A\ge m_B+m_C$, 第二括号与之共同确保实动量. 若 $C$ 无质量, $p_C=(m_A^2-m_B^2)/(2m_A)$. 这是静止 pion 衰变中中微子能量的来源.

无质量末态与 boost 满足 $E=\gamma E^*(1+\beta\cos\theta^*)$ 和
$\cos\theta=(\cos\theta^*+\beta)/(1+\beta\cos\theta^*)$. 消去 $\theta^*$ 得
\begin{equation}\label{eq:offaxis}
 E_\nu=\frac{E_\nu^*}{\gamma(1-\beta\cos\theta)}
 \simeq\frac{(1-m_\mu^2/m_\pi^2)E_\pi}{1+\gamma^2\theta^2}.
\end{equation}
后式要求 $\gamma\gg1$、$\theta\ll1$. 固定离轴角时, 中微子能量对 pion 能量存在极大值, 由 $\gamma\theta\simeq1$ 给出. 这不是说所有无质量末态角度都小于 $1/\gamma$: 反向发射仍能出现在很大实验室角度, 只是能量低.

\originalsection{Lorentz 不变相空间}
每个末态粒子的测度为 $\dd^3p/[(2\pi)^32E]$, 能动量守恒由四维 delta 实现:
\begin{equation}
 \dd\Phi_n=(2\pi)^4\delta^{(4)}(P-\sum_i p_i)\prod_i\frac{\dd^3p_i}{(2\pi)^32E_i}.
\end{equation}
二体质心相空间先用三维动量 $\delta^{(3)}(\mathbf p_B+\mathbf p_C)$ 置 $\mathbf p_C=-\mathbf p_B$，令 $p=|\mathbf p_B|$。剩下 $g(p)=\sqrt s-E_B(p)-E_C(p)$ 的能量 delta；在根 $p_f$ 上 $E_B+E_C=\sqrt s$，所以
\begin{equation}
\left|\frac{\dd g}{\dd p}\right|_{p_f}
=p_f\left(\frac1{E_B}+\frac1{E_C}\right)
=\frac{p_f\sqrt s}{E_BE_C},\qquad
\int_0^\infty\frac{p^2\dd p}{E_BE_C}\delta(g(p))
=\frac{p_f}{\sqrt s}.
\end{equation}
把 $\dd^3p=p^2\dd p\dd\Omega$ 与 $(2\pi)^4/(2\pi)^6$、两个 $1/(2E)$ 因子合并，才得到
\begin{equation}
 \frac{\dd\Phi_2}{\dd\Omega}=\frac{|\mathbf p_f|}{16\pi^2\sqrt s},\quad
 \frac{\dd\Gamma}{\dd\Omega}=\frac{|\mathcal M|^2|\mathbf p_f|}{32\pi^2m_A^2}.
\end{equation}
第二式另用了母粒子静止时 $\dd\Gamma=|\mathcal M|^2\dd\Phi_2/(2m_A)$；阈值 $p_f=0$ 处先取极限而非把雅可比分母单独置零。
相同末态粒子在完整相空间积分时另除 $n!$. 对无偏振初态要平均其自旋, 对未测量末态自旋则求和. 这是不同的操作: 末态不是随机制备, 所以不应平均掉其计数.

\originalsection{散射通量}
不变入射通量分母为 $4\sqrt{(p_a\cdot p_b)^2-m_a^2m_b^2}$. 在质心系等于 $4p_i\sqrt s$, 给
\begin{equation}\label{eq:cmcross}
 \frac{\dd\sigma}{\dd\Omega}=\frac1{64\pi^2s}\frac{p_f}{p_i}\overline{|\mathcal M|^2}.
\end{equation}
本式默认末态有标记或尚未加同粒子因子. 低速吸收反应常有 $\sigma\propto1/v$, 但 $\sigma v$ 可以有限. 发散的几何截面因子未必意味着物理反应率无限.
